\section{Introduction}

\begin{guideline}
L'introduction pose le fil directeur de l'ensemble du rapport. Elle doit
permettre à un lecteur de comprendre immédiatement : dans quel contexte
vous avez effectué ce stage, quel problème vous avez eu à adresser, et
comment vous allez le démontrer à travers ce rapport. Environ 1 à 2 pages.
\end{guideline}


\subsection{Contexte général}

\instruction{[Présentez le contexte dans lequel s'inscrit votre stage :
secteur d'activité, enjeux actuels du domaine, pourquoi ce sujet est
pertinent aujourd'hui. Ex. : montée en puissance des cybermenaces,
transformation digitale des organisations, adoption du cloud...]}


\subsection{Sujet et problématique}

\instruction{[Énoncez clairement le sujet de votre stage et la
problématique centrale que vous avez eu à adresser. La problématique
doit être formulée comme une question ou un défi concret. Ex. :
« Comment automatiser le déploiement et le durcissement d'infrastructures
Windows dans un environnement Azure en garantissant traçabilité et
sécurité de la chaîne d'approvisionnement logicielle ? »]}


\subsection{Fil directeur du rapport}

\instruction{[Expliquez brièvement comment votre rapport est structuré
et pourquoi dans cet ordre. Donnez au lecteur une « carte » de ce qu'il
va lire. Ex. : « Nous présenterons d'abord l'entreprise et le contexte
technique, avant de détailler les missions réalisées, les apprentissages
acquis et la valeur produite. »]}